\documentclass[conference]{IEEEtran}

\usepackage{amsmath}
\usepackage{booktabs}
\usepackage{graphicx}
\usepackage{siunitx}
\usepackage[hidelinks]{hyperref}

\title{Participant-Disjoint Evaluation of COVID-19 Respiratory-Audio Models}

\author{
\IEEEauthorblockN{Nishant Harkut and Ritu Tiwari}
\IEEEauthorblockA{Department of Information Technology\\
ABV-IIITM Gwalior, India}
}

\begin{document}
\maketitle

\begin{abstract}
COVID-19 audio screening models report AUROCs of 0.90--0.97. Han \emph{et al.} showed biased splits inflate AUROC from 0.71 to 0.90. Coppock \emph{et al.} found matching reduces audio AUROC from 0.85 to 0.62. We evaluate 2,088 Coswara participants under participant-disjoint splits with training-only feature ranking and validation-only model selection. The cough--speech fusion achieved test AUROC 0.895 (95\% CI 0.852--0.933, N=314). Speech alone achieved 0.888 (95\% CI 0.842--0.930). The fusion gain of 0.007 AUROC was not significant (DeLong $p=0.62$, 95\% CI for difference: $-0.021$ to $0.035$). Metadata alone achieved 0.964 (95\% CI 0.938--0.984). Shuffled-label control yielded 0.503. Results show strong internal discrimination but do not establish external validity.
\end{abstract}

\begin{IEEEkeywords}
COVID-19, respiratory audio, evaluation design, participant-disjoint validation
\end{IEEEkeywords}

\section{Introduction}

Audio-based COVID-19 screening reports AUROCs of 0.90--0.97 \cite{laguarta2020cough,bhattacharya2023coswara,aytekin2024hst}. Evaluation design affects these estimates: Han \emph{et al.} found biased participant allocation inflates AUROC from 0.71 to 0.90 \cite{han2022realistic}, Coppock \emph{et al.} found matching reduces audio AUROC from 0.85 to 0.62 while symptoms outperform audio \cite{coppock2024audio}, and Avila \emph{et al.} showed non-nested feature selection produces optimistic estimates \cite{avila2021feature}.

We evaluate Coswara under participant-disjoint splits with training-only feature ranking and validation-only model selection. Contributions:
\begin{enumerate}
\item Feature ranking uses training partition only (N=1,460), frozen for validation and test
\item Modality models and fusion selected by validation AUROC, test metrics excluded
\item Test results with bootstrap CIs and paired statistical tests
\item Metadata and shuffled-label baselines establish non-audio predictiveness
\end{enumerate}

The study evaluates internal performance. Temporal and cross-dataset validation require separate protocols.

\section{Related Work}

Coswara provides multimodal respiratory recordings with symptoms \cite{bhattacharya2023coswara}. COUGHVID offers large-scale cough recordings \cite{orlandic2021coughvid}. Chetupalli \emph{et al.} achieved AUROC 0.88 with subject-disjoint audio fusion \cite{chetupalli2023multimodal}. Truong \emph{et al.} reported AUROC $0.866\pm0.012$ with seven sound instances \cite{truong2024fair}. Aytekin \emph{et al.} proposed hierarchical spectrogram transformers for respiratory audio \cite{aytekin2024hst}.

Han \emph{et al.} demonstrated evaluation design changes AUROC from 0.71 to 0.90 \cite{han2022realistic}. Coppock \emph{et al.} showed matching reduces audio AUROC from 0.85 to 0.62, symptoms achieve 0.68 \cite{coppock2024audio}. Islam \emph{et al.} reported internal AUROC 0.92 but cross-dataset transfer of 0.53 \cite{islam2026robust}. These results indicate evaluation protocol determines reported performance.

\section{Methods}

\subsection{Problem Formulation}

Let $\mathcal{D} = \{(x_i, y_i)\}_{i=1}^N$ denote the dataset with features $x_i$ and labels $y_i \in \{0, 1\}$. We partition $\mathcal{D}$ into train $\mathcal{D}_{\text{train}}$, validation $\mathcal{D}_{\text{val}}$, and test $\mathcal{D}_{\text{test}}$ such that all recordings from participant $p$ remain in one partition:
\begin{equation}
\forall p: \quad \{x_i : i \in p\} \subset \mathcal{D}_k \text{ for exactly one } k \in \{\text{train, val, test}\}.
\end{equation}

\subsection{Dataset and Preprocessing}

The indexed Coswara release contained 2,746 participants. We excluded 632 with unresolved labels and 26 failing quality checks. The modeling cohort comprised 2,088 participants (Table~\ref{tab:cohort}).

\begin{table}[!b]
\caption{Participant-Disjoint Coswara Partitions}
\label{tab:cohort}
\centering
\footnotesize
\begin{tabular}{lrrrr}
\toprule
Partition & $N$ & Positive & Negative & Prevalence \\
\midrule
Train      & 1,460 & 476 & 984 & 32.6\% \\
Validation &   312 & 101 & 211 & 32.4\% \\
Test       &   316 & 103 & 213 & 32.6\% \\
\bottomrule
\end{tabular}
\end{table}

Audio was resampled to 16~kHz, silence-trimmed at 35~dB, and padded to $T$ seconds ($T=4$ for cough, $T=5$ for breathing and speech). Quality screening required minimum duration (0.5~s cough, 1.0~s breathing/speech), maximum silence ratio $\rho_s \leq 0.70$, and maximum clipping ratio $\rho_c \leq 0.01$.

\subsection{Feature Extraction}

We extracted three descriptor banks. Let $\mathbf{f}_{\text{ComParE}} \in \mathbb{R}^{6373}$ denote ComParE 2016 functionals \cite{schuller2016compare}, $\mathbf{f}_{\text{IS10}} \in \mathbb{R}^{1586}$ denote INTERSPEECH 2010 functionals \cite{schuller2010is10}, and $\mathbf{f}_{\text{proj}} \in \mathbb{R}^{2177}$ denote project acoustic features. The merged representation:
\begin{equation}
\mathbf{f} = [\mathbf{f}_{\text{ComParE}}^\top, \mathbf{f}_{\text{IS10}}^\top, \mathbf{f}_{\text{proj}}^\top]^\top \in \mathbb{R}^{10140}.
\end{equation}

\subsection{Feature Ranking}

LightGBM gain importance $g_j$ for feature $j$ was computed \emph{only} on $\mathcal{D}_{\text{train}}$. For each modality $m$, we normalized:
\begin{equation}
s_{j,m} = \frac{g_{j,m}}{\max_k g_{k,m}},
\end{equation}
then averaged across modalities:
\begin{equation}
s_j = \frac{1}{|\mathcal{M}|} \sum_{m \in \mathcal{M}} s_{j,m}.
\end{equation}

The 800 highest-scoring features were frozen. The same feature indices were used in $\mathcal{D}_{\text{val}}$ and $\mathcal{D}_{\text{test}}$ without recomputing $s_j$.

\subsection{Modality Models}

We evaluated four classifiers: LightGBM \cite{ke2017lightgbm}, XGBoost \cite{chen2016xgboost}, CatBoost \cite{prokhorenkova2018catboost}, and RBF-SVC \cite{cortes1995svm}. Let $f_\theta(x)$ denote the probability output for input $x$. Recording-level probabilities were aggregated:
\begin{equation}
p_m(x) = \frac{1}{|R_p|} \sum_{r \in R_p} f_\theta(x_r),
\end{equation}
where $R_p$ denotes recordings from participant $p$ for modality $m$.

For each modality, validation AUROC selected either one classifier or the mean of top performers.

\subsection{Fusion}

We evaluated three fusion rules on $\mathcal{D}_{\text{val}}$:

\textbf{Uniform mean:}
\begin{equation}
p_{\text{fusion}}(x) = \frac{1}{|\mathcal{M}_x|} \sum_{m \in \mathcal{M}_x} p_m(x).
\end{equation}

\textbf{Validation-weighted:} Weights $w_m$ derived from validation AUPRC:
\begin{equation}
w_m = \frac{\max(\text{AUPRC}_m - 0.5, 0.01)}{\sum_j \max(\text{AUPRC}_j - 0.5, 0.01)}.
\end{equation}

\textbf{Logistic stack:} Class-balanced L2 logistic regression:
\begin{equation}
p_{\text{stack}}(x) = \sigma\left(\beta_0 + \sum_m \beta_m p_m(x)\right).
\end{equation}

The primary rule was uniform averaging. Validation AUROC selected \emph{which modalities} to include. Test metrics were excluded from selection.

\subsection{Threshold Selection}

Balanced-accuracy threshold $\tau$ was selected on $\mathcal{D}_{\text{val}}$:
\begin{equation}
\tau = \arg\max_t \frac{1}{2}\left(\text{TPR}(t) + \text{TNR}(t)\right).
\end{equation}

\subsection{Baselines}

\textbf{Metadata model:} LightGBM trained on demographic and recording-protocol variables.

\textbf{Shuffled-label control:} Complete pipeline rerun with $\tilde{y}_i = y_{\pi(i)}$ for random permutation $\pi$.

\subsection{Evaluation Metrics}

Primary endpoint was participant-level AUROC. For binary classifier with true positive rate $\text{TPR}(t)$ and false positive rate $\text{FPR}(t)$ at threshold $t$:
\begin{equation}
\text{AUROC} = \int_0^1 \text{TPR}(t) \, d(\text{FPR}(t)).
\end{equation}

We also reported AUPRC, balanced accuracy, and F1:
\begin{align}
\text{Balanced Acc.} &= \frac{1}{2}(\text{TPR} + \text{TNR}), \\
\text{F1} &= \frac{2 \cdot \text{TP}}{2 \cdot \text{TP} + \text{FP} + \text{FN}}.
\end{align}

\subsection{Uncertainty Quantification}

Nonparametric bootstrap resampled test participants $B=1000$ times. For AUROC $\hat{A}$, the 95\% CI:
\begin{equation}
[\hat{A}_{0.025}, \hat{A}_{0.975}],
\end{equation}
where $\hat{A}_\alpha$ is the $\alpha$ quantile of bootstrap distribution.

Paired DeLong test compared AUROCs where same participants had predictions from both models \cite{delong1988auc}.

\section{Results}

\subsection{Validation-Only Selection}

Validation AUROC selected cough+speech (0.842) over cough+breath+speech (0.841). Validation AUROC selected LightGBM for speech and four-model mean for cough and breathing.

\subsection{Test Performance}

Table~\ref{tab:results} shows test performance. Cough+speech fusion achieved AUROC 0.895 (95\% CI 0.852--0.933, N=314). Speech alone achieved 0.888 (95\% CI 0.842--0.930). The fusion gain of 0.007 AUROC was not statistically significant (DeLong test, difference 0.007, 95\% CI $-0.021$ to $0.035$, $p=0.62$).

\begin{table}[!t]
\caption{Test Results with 95\% Bootstrap Confidence Intervals}
\label{tab:results}
\centering
\footnotesize
\begin{tabular}{lrrrr}
\toprule
System & $N$ & AUROC (95\% CI) & AUPRC & F1 \\
\midrule
Breathing, mean-4  & 312 & 0.828 (0.775--0.877) & 0.734 & 0.654 \\
Cough, mean-4      & 316 & 0.862 (0.812--0.908) & 0.804 & 0.705 \\
Speech, LightGBM   & 314 & 0.888 (0.842--0.930) & 0.833 & 0.740 \\
\midrule
\textbf{Cough+speech} & \textbf{314} & \textbf{0.895 (0.852--0.933)} & \textbf{0.862} & \textbf{0.729} \\
\bottomrule
\end{tabular}
\end{table}

Fusion F1 (0.729) was lower than speech F1 (0.740). Balanced accuracy was 0.808 for fusion vs. 0.809 for speech.

\subsection{Baseline Comparisons}

Metadata model achieved AUROC 0.964 (95\% CI 0.938--0.984). Shuffled-label control yielded mean AUROC 0.503.

\subsection{Literature Context}

Table~\ref{tab:context} compares with prior Coswara studies. These are not head-to-head comparisons: cohort snapshots, labels, and partitions differ.

\begin{table}[!t]
\caption{Protocol-Aware Literature Comparison}
\label{tab:context}
\centering
\footnotesize
\begin{tabular}{llc}
\toprule
Study & Protocol & AUROC \\
\midrule
Bhattacharya \cite{bhattacharya2023coswara} & Internal, 9 sounds + symptoms & 0.915 \\
Chetupalli \cite{chetupalli2023multimodal} & Subject-disjoint, audio fusion & 0.880 \\
Truong \cite{truong2024fair} & Fixed test, 7 sounds & $0.866\pm0.012$ \\
\midrule
This work & Participant-disjoint, 2 sounds & 0.895 \\
\bottomrule
\end{tabular}
\end{table}

\section{Discussion}

Participant-disjoint evaluation yielded AUROC 0.895 (95\% CI 0.852--0.933), within the range of prior internal estimates (0.88--0.92). The fusion gain over speech (0.007 AUROC) was not significant ($p=0.62$). F1 decreased from 0.740 (speech) to 0.729 (fusion).

Metadata achieved AUROC 0.964, higher than audio (0.895), consistent with Han and Coppock's finding that non-audio variables are predictive \cite{han2022realistic,coppock2024audio}. Shuffled-label AUROC 0.503 confirms absence of leakage.

\textbf{Limitations:} First, labels were self-reported, not PCR-confirmed. Second, validation served multiple selection roles; nested evaluation would be more conservative. Third, only 314 participants had both cough and speech predictions. Fourth, this evaluates internal performance only; external and temporal validation require separate protocols. Fifth, breathing was excluded from the primary model.

\section{Conclusion}

Under participant-disjoint evaluation with training-only feature ranking and validation-only model selection, a Coswara audio model achieved AUROC 0.895 (95\% CI 0.852--0.933). Fusion did not significantly outperform speech alone ($p=0.62$). Metadata outperformed audio (AUROC 0.964 vs. 0.895). The estimate applies only to the test cohort and does not establish external validity.

\bibliographystyle{IEEEtran}
\bibliography{references}

\end{document}